\documentclass[11pt]{article}
\usepackage[T1]{fontenc}
\usepackage{lmodern,amsmath,amssymb,amsthm,mathtools}
\usepackage[margin=1.05in]{geometry}
\usepackage[hidelinks]{hyperref}
\newtheorem{theorem}{Theorem}[section]
\newtheorem{lemma}[theorem]{Lemma}
\newtheorem{proposition}[theorem]{Proposition}
\newtheorem{corollary}[theorem]{Corollary}
\newcommand{\R}{\mathbb R}\newcommand{\E}{\mathbb E}
\newcommand{\Hh}{\mathcal H}\newcommand{\code}[1]{{\small\nolinkurl{#1}}}
\title{Gaussian clock-domain obstructions\protect\\
for reflected-clock Brownian motion}
\author{Project Stardust}
\date{Reviewed editorial draft, October 29, 2026}
\begin{document}
\maketitle
\begin{abstract}
Let $B,W$ be independent Brownian motions and
$Z_t=B(|W_t|)$. At each positive time the actual point
observable belongs to the scenery Gaussian Sobolev
domain, with an explicit finite-energy gradient, but
does not belong to the original clock Gaussian Sobolev
domain. The bounded smooth observable $\sin Z_t$ has the
same domain separation, so unbounded point growth is not
the cause. The obstruction is proved in the scenery-valued
Hilbert space: its clock section has squared increment
norm proportional to distance, incompatible with a
common locally finite-energy weak derivative. A
conditional contraction of the full clock gradient
would recover precisely such a derivative, so separate
scalar identities cannot evade the obstruction.
An exact correlated-clock noise formula further gives a sharp
weighted integral-energy cutoff and divergence of the original
square-root-weighted clock-chaos energy.
For a closed densely defined Hilbert operator, strong
value approximation of a point outside its domain forces
all derivative norms to escape to infinity, while the
embedded derivatives converge in the dual of its
adjoint graph with an explicit error bound. Applied to
the ambient clock operator, this gives a nonrepresentable
graph-dual derivative of $Z_t$. The conclusions concern
the stated original domains and do not supply an
observed-filtration stochastic integral or a physical-time
It\^o formula. The proofs reconstruct frozen Lean results
using classical Gaussian and closed-operator methods;
novelty remains unresolved.
\end{abstract}

\section{The actual process and the two Gaussian domains}
On a product of independent Wiener probability spaces
$(\Omega_B,\mathbb P_B)$ and $(\Omega_W,\mathbb P_W)$,
let $B_s,W_t$, $s,t\geq0$, be standard Brownian motions
starting at zero. We study
\begin{equation}\label{eq:process}
 Z_t(\omega_B,\omega_W)=B_{|W_t(\omega_W)|}(\omega_B).
\end{equation}
This is the simple Brownian-time Brownian motion of
Allouba--Zheng~\cite{alloubazheng}. The name
``reflected-clock'' specifies that the same scenery
$B$ is used at every visit to a given nonnegative
clock level. It does not mean an excursion-by-excursion
replacement of scenery or a two-sided scenery evaluated
at the signed clock.

Set
\[
 E=L^2(\mathbb P_B\otimes\mathbb P_W),\qquad
 H_B=H_W=L^2(\R_+;ds).
\]
The isonormal coordinates of $B$ and $W$ identify their
Cameron--Martin derivative directions with $H_B$ and
$H_W$. Write
\[
 D_B:D(D_B)\subset E\longrightarrow L^2(\Omega_B\times\Omega_W;H_B),
 \qquad
 D_W:D(D_W)\subset E\longrightarrow L^2(\Omega_B\times\Omega_W;H_W)
\]
for the original closed partial Gaussian derivatives.
The other source is held as an $L^2$ parameter in each
case; these are not derivatives in physical time $t$.
Section~\ref{sec:gaussian} specifies their closure and
fiberwise domains precisely.

\begin{theorem}\label{thm:point}
For every $T>0$, $Z_T\in D(D_B)$ and
\begin{equation}\label{eq:sourcegradient}
 D_B Z_T=\mathbf1_{[0,|W_T|]}\quad\hbox{in }H_B,
 \qquad
 \|D_B Z_T\|_2^2=\|Z_T\|_2^2=\sqrt{2T/\pi}.
\end{equation}
In contrast, $Z_T\notin D(D_W)$. Consequently the
point is outside any joint derivative whose domain is
contained in $D(D_B)\cap D(D_W)$, including the
original joint derivative restricted to IBM-observable
$L^2$ classes in the mapped development. At $T=0$ both
derivatives of $Z_0=0$ exist and are zero.
\end{theorem}

\begin{theorem}\label{thm:bounded}
For every $T>0$, $S_T=\sin Z_T$ is a bounded actual
observable in $D(D_B)$, with
\[
 D_BS_T=\cos Z_T\,\mathbf1_{[0,|W_T|]},\qquad
 \|D_BS_T\|_2^2\leq\sqrt{2T/\pi}.
\]
Nevertheless $S_T\notin D(D_W)$ and it is outside the
original joint intrinsic domain. Thus that domain does
not contain all bounded smooth functions of actual point
observations.
\end{theorem}

We will also prove the following quantitative form of
the distinction between the original graph and a weak
test-dual extension.

\begin{theorem}\label{thm:dualpoint}
Let $\mathcal T=D(D_W^*)$ with norm
$\|v\|_{\mathcal T}^2=\|v\|_2^2+\|D_W^*v\|_2^2$.
For either $f=Z_T$ or $f=S_T$ there is a bounded clock graph-dual derivative
\[
 d_W f(v)=\E[f D_W^*v],\qquad
 \|d_W f\|_{\mathcal T^*}\leq\|f\|_2.
\]
For $T>0$ no $V\in L^2(\Omega_B\times\Omega_W;H_W)$
represents it as $\E\langle V,v\rangle_{H_W}$ for all
$v\in\mathcal T$. Whenever $u_n\in D(D_W)$ and
$u_n\to f$ in $E$,
\begin{align}
 \|D_Wu_n\|_2&\longrightarrow\infty,\label{eq:escape}\\
 \|jD_Wu_n-d_W f\|_{\mathcal T^*}
      &\leq\|u_n-f\|_2\longrightarrow0,\label{eq:dualerror}
\end{align}
where $jV(v)=\E\langle V,v\rangle_{H_W}$.
\end{theorem}

This theorem concerns the ambient original clock
operator. Its tests can depend on both sources and are
not identified with temporal predictable integrands or
with the natural filtration of $Z$. The native
point-domain obstruction in Theorem~\ref{thm:point}
is a separate consequence of domain inclusion.

\section{Closed Gaussian derivatives and their fiberwise lift}
\label{sec:gaussian}
Here are the Gaussian foundations used in the proof.
For a real separable Hilbert space $H$, let $X(h)$ be an
isonormal Gaussian family, with
$\E X(h)X(k)=\langle h,k\rangle_H$.
On finite polynomial cylinders define
\begin{equation}\label{eq:cylinder}
 D\,p(X(h_1),\ldots,X(h_n))
   =\sum_{i=1}^n(\partial_i p)(X(h_1),\ldots,X(h_n))h_i.
\end{equation}
One may use orthonormal $h_i$ after a finite linear
change of coordinates. Polynomial cylinders are dense
in Gaussian $L^2$. For completeness, in finite dimension
if $f$ is orthogonal to every polynomial, the entire
transform $\int f(x)e^{z\cdot x}\,d\gamma_n(x)$ has
all Taylor coefficients zero. Cauchy--Schwarz and
Gaussian exponential moments justify its entireness and
termwise differentiation. Its restriction to imaginary
$z$ is the Fourier transform of $f\gamma_n$, so Fourier
uniqueness gives $f=0$. In countably many coordinates,
conditional expectations onto the first $n$ coordinates
converge in $L^2$, reducing density to the finite case.

Gaussian integration by parts gives, for polynomial
cylinders $p,a$ and a finite coordinate vector $h$,
\begin{equation}\label{eq:stein}
 \E[\langle Dp,h\rangle a]
 =\E[p\{X(h)a-D_h a\}].
\end{equation}
This proves closability of \eqref{eq:cylinder}.
Indeed if $p_n\to0$ and $Dp_n\to V$ in their $L^2$
spaces, the right side tends to zero for each fixed
$a,h$. Such scalar-polynomial vector tests are dense
in $L^2(\Omega;H)$, so $V=0$.
The original Gaussian derivative is the closure of this
graph. Its domain is dense by the preceding density
argument. Smooth compactly supported coordinate tests
also satisfy \eqref{eq:stein}, by ordinary finite-dimensional
integration by parts and passage through the closed graph.
These standard source-space constructions are the setting
of Gaussian Malliavin calculus; see~\cite{sottinen} for
the corresponding Gaussian-process conventions.

For an independent parameter probability space
$(S,\rho)$, the partial derivative on
$L^2(S;L^2(\Omega))$ has the exact domain
\begin{equation}\label{eq:fiber}
\begin{split}
 D(\widetilde D)=\{f: &\text{there exists }V\in L^2(S;L^2(\Omega;H))\text{ with}\\
 &(f(s),V(s))\in\operatorname{graph}D
           \text{ for }\rho\text{-almost every }s\}.
\end{split}
\end{equation}
The output $V$ is unique, since the graph has a unique
value in each fiber. This is the actual Bochner graph
lift, not a definition using independently chosen
scalar derivatives.

The lift is closed: from a convergent sequence of both
graph coordinates, take a common subsequence converging
almost everywhere in both fiber norms, and use closedness
of the original graph in each fiber. It is densely
defined because simple $S$-functions with values in
$D(D)$ are dense in the Bochner $L^2$ space.
It is also the closure of tensor cylinders. To see the
nontrivial inclusion, a graph pair in \eqref{eq:fiber}
is a strongly measurable, square-integrable function
into the closed Hilbert graph of $D$. Approximate it
by simple graph-valued functions, and approximate each
of their finitely many values by a polynomial graph
pair of $D$. This yields tensor-cylinder approximation
in both graph coordinates. No measurable selection of
derivative values is required: the given pair already
supplies the measurable output.

Taking $S=\Omega_W$ or $\Omega_B$ defines $D_B$ and
$D_W$ as used above. Identifying product scalar $L^2$
with either Bochner order is an isometric Fubini
identification; changing the order does not change the
product law or the graph condition.

\section{Finite scenery energy at a random clock level}
Let $I_B:H_B\to L^2(\mathbb P_B)$ be the Wiener integral,
so $B_s=I_B(\mathbf1_{[0,s]})$.
Fix $s\geq0$ and an orthonormal basis $(e_j)$ of $H_B$.
The partial sums
\[
 B_s^{(N)}=\sum_{j=1}^N
     \langle\mathbf1_{[0,s]},e_j\rangle I_B(e_j)
\]
are polynomial cylinders, with gradients equal to the
orthogonal projections of $\mathbf1_{[0,s]}$.
Wiener isometry and Parseval show convergence of both
values and gradients. Thus $B_s$ is in the original
scenery domain, with derivative $\mathbf1_{[0,s]}$.

Now put $s=|W_T|$. The functions of the clock multiplying
$I_B(e_j)$ are square integrable, because
\[
 |\langle\mathbf1_{[0,|W_T|]},e_j\rangle|^2\leq|W_T|.
\]
They are allowed $L^2$ parameters in the scenery lift.
The random projected gradients converge to
$\mathbf1_{[0,|W_T|]}$ in the full product $L^2$ field
norm, by pointwise Parseval and domination by $|W_T|$.
The same is true of the random partial-sum values by
conditioning on $W$. Closedness proves the scenery
part of Theorem~\ref{thm:point}. Conditioning also gives
\[
 \E Z_T^2=\E|W_T|=\sqrt{2T/\pi};
\]
the last equality is the elementary absolute first
moment of a centered Gaussian of variance $T$.
This establishes every identity in
\eqref{eq:sourcegradient}.

\section{A Hilbert-valued section with Brownian metric}
Set $\Hh=L^2(\mathbb P_B)$ and, for $c>0$, define
\begin{equation}\label{eq:section}
 F_c(x)=B_{c|x|}\in\Hh,\qquad x\in\R.
\end{equation}
It is a continuous Hilbert-valued curve. For nonnegative
$x,y$, Wiener isometry gives the exact metric
\begin{equation}\label{eq:metric}
 \|F_c(x)-F_c(y)\|_{\Hh}^2=c|x-y|.
\end{equation}
In particular $F_c\in L^2(\gamma;\Hh)$ for standard
one-dimensional Gaussian measure $\gamma$, since
$\|F_c(x)\|_{\Hh}^2=c|x|$.

\begin{lemma}\label{lem:energy}
If a continuous Banach-valued curve $F$ on $[a,b]$,
$a<b$, satisfies
$\|F(x)-F(a)\|^2\geq c(x-a)$ for $c>0$,
then it has no representation
\[
 F(x)=v+\int_a^xg(s)\,ds
\]
almost everywhere on $[a,b]$ with $g\in L^2([a,b])$
and a constant vector $v$.
\end{lemma}
\begin{proof}
The primitive is continuous, so almost-everywhere
equality with the continuous $F$ implies equality
everywhere. At $a$ this gives $v=F(a)$.
Cauchy--Schwarz for the Bochner integral gives, for $x>a$,
\[
 c(x-a)\leq\|F(x)-F(a)\|^2
 \leq(x-a)\int_a^x\|g(s)\|^2\,ds.
\]
After division by $x-a$, the right side tends to zero
as $x\downarrow a$, a contradiction.
\end{proof}

Gaussian weighting does not remove this local
obstruction: on a compact interval its density is
bounded above and bounded away from zero, so Gaussian
$L^2$ derivative energy implies Lebesgue $L^2$ energy
there.

\begin{lemma}\label{lem:weakvector}
Let $\Hh$ be a real separable Hilbert space and
$F,g\in L^2(\gamma;\Hh)$. Suppose every scalar pair
$\langle e,F\rangle,\langle e,g\rangle$, $e\in\Hh$,
satisfies the closed one-dimensional Gaussian derivative
identity. Then on every finite interval $[a,b]$,
\[
 F(x)=v+\int_a^xg(s)\,ds
\]
almost everywhere, for some $v\in\Hh$.
\end{lemma}
\begin{proof}
The scalar graph identity implies
\[
 \int\langle e,F(x)\rangle\{xp(x)-p'(x)\}\,d\gamma(x)
       =\int\langle e,g(x)\rangle p(x)\,d\gamma(x)
\]
for compactly supported smooth $p$. If $\varphi$ is a
compactly supported test function in $(a,b)$, take
$p=\varphi/\rho$, where $\rho$ is the Gaussian density.
Since $\rho'=-x\rho$, $xp-p'=-\varphi'/\rho$.
Thus $\langle e,g\rangle$ is the ordinary weak derivative
of $\langle e,F\rangle$ on $(a,b)$.
Both are locally $L^2$, so the one-dimensional weak
fundamental theorem gives
$\langle e,F(x)\rangle=c_e+
\langle e,\int_a^xg(s)\,ds\rangle$ almost everywhere.
This elementary theorem follows, for example, by
subtracting the displayed primitive and observing that
a locally integrable distribution with zero derivative
is constant; convolution with compact smooth kernels
proves that assertion directly.

Choose a countable dense sequence $(e_j)$ in $\Hh$.
Intersect their full-measure equality sets and choose
one point $x_0$ in the intersection. The vector
$v=F(x_0)-\int_a^{x_0}g(s)\,ds$ realizes every constant
$c_{e_j}$. At any other point in the intersection,
the inner products with all $e_j$ of
$F(x)-v-\int_a^xg$ vanish. Density makes the vector
itself zero, proving the statement.
\end{proof}

Combining \eqref{eq:metric} on $[0,1]$ with the two
lemmas shows that $F_c$ has no single
$g\in L^2(\gamma;\Hh)$ realizing all scalar closed
Gaussian derivative pairs. This is a common-field
obstruction. It does not assert that every individual
scalar weak derivative fails to exist.

\section{Recovering a common field from the full clock graph}
We now connect the section obstruction to the actual
infinite-dimensional clock operator. Let
\[
 h_T=T^{-1/2}\mathbf1_{[0,T]}\in H_W,\qquad
 G=W(h_T)=W_T/\sqrt T.
\]
Then $\|h_T\|=1$ and $G$ is standard Gaussian.
Complete $h_T$ to an orthonormal basis of $H_W$.
The clock Gaussian space decomposes into the coordinate
$G$ and an independent residual Gaussian family.
Under this decomposition,
\begin{equation}\label{eq:pullback}
 Z_T=F_{\sqrt T}(G)
\end{equation}
as an actual product $L^2$ class, with $F$ valued in
the scenery Hilbert space.

Suppose, for contradiction, that $Z_T\in D(D_W)$,
with full clock gradient $V$. Contract $V$ with $h_T$
and average only over the residual clock coordinates:
\begin{equation}\label{eq:recover}
 g(x)(\omega_B)=
 \E_{\rm residual}\langle V(\omega_B,x,\cdot),h_T\rangle_{H_W}.
\end{equation}
Hilbert-valued conditional integration and Jensen give
\[
 g\in L^2(\gamma;\Hh),\qquad
 \|g\|_{L^2(\gamma;\Hh)}\leq\|V\|_2.
\]
This is one common scenery-valued field, rather than
separately chosen scalar derivatives.

For $e\in\Hh$ and $p\in C_c^\infty(\R)$, the
clock test field $e(\omega_B)p(G)h_T$ has divergence
\[
 e(\omega_B)\{G p(G)-p'(G)\}.
\]
The formula follows on tensor polynomial cylinders
from Gaussian integration by parts. Its test field
and displayed divergence are in the relevant $L^2$
spaces; passage through the closed clock graph extends
the identity to $Z_T,V$. Using \eqref{eq:pullback}
and conditional integration in \eqref{eq:recover}, it
becomes
\[
 \int\langle e,F_{\sqrt T}(x)\rangle
          \{xp(x)-p'(x)\}\,d\gamma(x)
 =\int\langle e,g(x)\rangle p(x)\,d\gamma(x).
\]
The proof of Lemma~\ref{lem:weakvector} therefore
supplies a common vector primitive for $F_{\sqrt T}$
on $[0,1]$. Lemma~\ref{lem:energy} excludes it by
\eqref{eq:metric}. This proves $Z_T\notin D(D_W)$.

The original joint intrinsic derivative in the mapped
library is defined on the intersection of the scenery
and clock domains, restricted to the closed subspace
of actual IBM-observable $L^2$ classes. The point
$Z_T$ is observable by definition but fails the clock
domain test. Thus the joint-domain conclusion requires
no assertion about a new derivative domain or a new
filtration. This completes Theorem~\ref{thm:point}.

\section{A bounded smooth observable with the same obstruction}
The preceding argument is not restricted to unbounded
point values. We prove Theorem~\ref{thm:bounded} using
an exact first-chaos projection.
Let $\Pi$ be the orthogonal projection in
$\Hh=L^2(\mathbb P_B)$ onto the closed first Gaussian
chaos $I_B(H_B)$. For $s\geq0$,
\begin{equation}\label{eq:sineprojection}
 \Pi\sin B_s=e^{-s/2}B_s.
\end{equation}
Indeed, for $h\in H_B$, joint Gaussian regression and
one-dimensional integration by parts give
\[
 \E[I_B(h)\sin B_s]
 =\langle h,\mathbf1_{[0,s]}\rangle\E\cos B_s
 =e^{-s/2}\langle h,\mathbf1_{[0,s]}\rangle.
\]
For $s>0$, regression here follows simply by splitting
$h$ into its projection onto $\mathbf1_{[0,s]}$ and
its orthogonal complement; the latter Gaussian is
independent and centered. At $s=0$ the identity is zero.
The last expression is the inner product with
$e^{-s/2}B_s$, proving \eqref{eq:sineprojection}.

For $c>0$ set $S_c(x)=\sin B_{c|x|}$ as a
scenery-valued Hilbert class. It is continuous and has
norm at most one. Its projected curve is
\[
 \Pi S_c(x)=e^{-c|x|/2}B_{c|x|},\qquad \Pi S_c(0)=0.
\]
On $0\leq x\leq1$, the exact source variance implies
\begin{equation}\label{eq:sinelower}
 \|\Pi S_c(x)-\Pi S_c(0)\|_{\Hh}^2
       =e^{-cx}cx\geq ce^{-c}x.
\end{equation}
If $S_c$ had a common finite-energy vector primitive,
applying the bounded operator $\Pi$ to its Bochner
integral would produce a primitive of $\Pi S_c$ with
no larger derivative energy. Lemma~\ref{lem:energy}
excludes this by \eqref{eq:sinelower}.
The recovery argument of the preceding section, with
$F_c$ replaced by $S_c$, would recover just such a
primitive from a full clock gradient of $\sin Z_T$.
Thus $\sin Z_T\notin D(D_W)$ at every positive time,
and domain inclusion gives the joint intrinsic exclusion.

For the positive scenery statement, at each deterministic
$s$ the Taylor polynomials of $\sin B_s$ and their
derivatives converge in Gaussian $L^2$, dominated by
a constant times $e^{|B_s|}$. Their gradients converge
to $\cos B_s\,\mathbf1_{[0,s]}$. Closedness gives that
actual scenery graph pair. At the random level
$s=|W_T|$ its field norm squared is bounded by $s$,
and its value is bounded by one. Measurability follows
from continuity of $B$ and the interval indicator map.
The exact fiber domain \eqref{eq:fiber} therefore gives
the scenery graph of $S_T$ and its norm bound in
Theorem~\ref{thm:bounded}.
This is not an attempted clock chain rule applied to
an input outside the clock domain.

\section{Exact clock-noise energy and a critical sector obstruction}
\label{sec:noise}
The clock section also gives quantitative information below ordinary
first-order differentiability. We keep the statements in terms of an
explicit coupling and the original orthogonal clock-chaos projections.
This avoids assuming an identification with a fractional operator domain.

Fix $T>0$, put $c=\sqrt T$ and $\kappa=\sqrt{2/\pi}$, and let
\[
 G=W_T/\sqrt T,\qquad G'=W_{T+1}-W_T.
\]
The variables $G,G'$ are independent standard normals, independent of
the scenery. For $0\leq h\leq1$, define
\[
 G_h=(1-h)G+\sqrt{h(2-h)}G',\qquad Z_T^{(h)}=B_{c|G_h|}.
\]
Then $G_h$ is standard normal and $\E GG_h=1-h$. The construction
uses the future clock increment $G'$; it is a Gaussian noise probe and
is not an adapted physical-time driver. The same full random scenery
is retained in both $Z_T$ and $Z_T^{(h)}$.

\begin{theorem}[Exact noise formula]\label{thm:noise}
For $0\leq h\leq1$,
\begin{equation}\label{eq:noiseexact}
 \|Z_T-Z_T^{(h)}\|_E^2
  =c\kappa r(h),\qquad
 r(h)=\sqrt{4-2h}+\sqrt{2h}-2.
\end{equation}
One has
\begin{equation}\label{eq:noisebounds}
 \frac{\sqrt{2h}}4\leq r(h)\leq\sqrt{2h},\qquad
 \lim_{h\downarrow0}\frac{r(h)}{\sqrt h}=\sqrt2.
\end{equation}
In particular, for any real $a$,
\begin{equation}\label{eq:noiseintegral}
 \int_0^1\frac{\|Z_T-Z_T^{(h)}\|_E^2}{h^{1+a}}\,dh<\infty
 \quad\Longleftrightarrow\quad a<\tfrac12.
\end{equation}
\end{theorem}
\begin{proof}
Conditioning on the clock and using the exact scenery increment metric
gives the squared norm $c\E\bigl||G|-|G_h|\bigr|$.
For arbitrary real $x,y$, checking their signs and which has larger
absolute value gives
\[
 \bigl||x|-|y|\bigr|=|x+y|+|x-y|-|x|-|y|.
\]
The centered Gaussian variables $G+G_h$ and $G-G_h$ have variances
$4-2h$ and $2h$, respectively. Since a centered normal with variance
$v$ has first absolute moment $\kappa\sqrt v$, this proves
\eqref{eq:noiseexact}.

For $h>0$, rationalization gives
\[
 r(h)=\sqrt{2h}-\frac{2h}{2+\sqrt{4-2h}},\qquad
 \frac{r(h)}{\sqrt h}
   =\sqrt2-\frac{2\sqrt h}{2+\sqrt{4-2h}}.
\]
The limit and upper bound follow. Moreover
\[
 \frac{r(h)}{\sqrt{2h}}
  =1-\frac{\sqrt{2h}}{2+\sqrt{4-2h}}
  \geq1-\frac{\sqrt2}{2}\geq\frac14
 \quad(0<h\leq1).
\]
The case $h=0$ is immediate. As $c\kappa>0$, the integrand
in~\eqref{eq:noiseintegral} is bounded above and below by positive
constant multiples of $h^{-a-1/2}$. Its integral at zero is finite
exactly when $-a-1/2>-1$. There is no other endpoint singularity.
\end{proof}

Let $\Pi_n^W$ be the original orthogonal projection onto the $n$th
homogeneous chaos of the clock Gaussian source, acting as the identity
on all scenery randomness. Put
\[
 E_n(T)=\|\Pi_n^W Z_T\|_E^2,\qquad n\geq0.
\]
These are energies of the actual endpoint, not of its scenery average.
They are nonnegative and sum to $\|Z_T\|_E^2=c\kappa$.

\begin{lemma}[Clock-sector noise identity]\label{lem:noiseseries}
For $0\leq h\leq1$,
\begin{equation}\label{eq:noiseseries}
 2\sum_{n\geq0}(1-(1-h)^n)E_n(T)=c\kappa r(h).
\end{equation}
Every odd clock projection of $Z_T$ vanishes.
\end{lemma}
\begin{proof}
The section $\Phi(x)=B_{c|x|}$ belongs to
$L^2(\gamma;L^2(\mathbb P_B))$, where $\gamma$ is standard Gaussian
measure. Its Hilbert-valued Hermite expansion is
\[
 \Phi(x)=\sum_{n\geq0}b_n\frac{\operatorname{He}_n(x)}{\sqrt{n!}},
 \qquad \sum_n\|b_n\|_{L^2(\mathbb P_B)}^2=c\kappa.
\]
Completeness follows by applying scalar Hermite completeness to the
coordinates of a countable Hilbert basis of $L^2(\mathbb P_B)$;
Parseval and nonnegative summation give the asserted expansion and
energy equality. A scalar function of the single normalized clock
direction $G$ has its $n$th Hermite component in the original full
clock chaos of order $n$. Finite Hermite expansions and continuity
of the orthogonal projections therefore identify
$E_n(T)=\|b_n\|^2$.

For standard normals $X,Y$ of correlation $\rho$, their exponential
generating functions give
\[
 \E[e^{sX-s^2/2}e^{tY-t^2/2}]=e^{\rho st},\qquad
 \E[\operatorname{He}_n(X)\operatorname{He}_m(Y)]
   =\mathbf1_{n=m}n!\rho^n.
\]
Differentiation at $s=t=0$ is justified by Gaussian exponential
integrability. For finite Hilbert-valued Hermite sums it follows that
the squared difference at $G,G_h$ is
$2\sum_n(1-(1-h)^n)\|b_n\|^2$.
Both pullbacks from $L^2(\gamma;L^2(\mathbb P_B))$ into the joint
space are isometries, so the finite sums converge in both pullback
norms. The series has nonnegative terms bounded by $2E_n(T)$.
Passing to the limit and using~\eqref{eq:noiseexact} proves
\eqref{eq:noiseseries}. Finally $\Phi(-x)=\Phi(x)$ whereas odd
Hermite polynomials change sign. Thus $b_n=0$ for odd $n$, and
the corresponding original clock projections vanish.
\end{proof}

\begin{theorem}[Critical weighted energy fails]\label{thm:critical}
For each $T>0$,
\begin{equation}\label{eq:criticalsqrt}
 \sum_{n\geq1}\sqrt n\,E_n(T)=\infty.
\end{equation}
Consequently, for every $a\geq\tfrac12$,
\[
 \sum_{n\geq0}(1+n)^aE_n(T)=\infty.
\]
Beyond every finite cutoff there is a nonzero even clock sector of
$Z_T$.
\end{theorem}
\begin{proof}
For $n\geq0$ and $0<h\leq1$, Bernoulli's inequality gives
\[
 0\leq1-(1-h)^n\leq\min(1,nh)\leq\sqrt{nh}.
\]
For each fixed $n$ the same bound by $nh$ gives
$(1-(1-h)^n)/\sqrt h\leq n\sqrt h\to0$.
If the sum in~\eqref{eq:criticalsqrt} were finite, dominated
convergence for counting measure, with dominating sequence
$2\sqrt n E_n(T)$, would imply
\[
 \lim_{h\downarrow0}
 2\sum_{n\geq0}\frac{1-(1-h)^n}{\sqrt h}E_n(T)=0.
\]
But~\eqref{eq:noiseseries} and~\eqref{eq:noisebounds} give the
strictly positive limit $c\kappa\sqrt2$, a contradiction.
For $a\geq1/2$, $(1+n)^a\geq\sqrt n$, proving the second
assertion. If all even projections beyond some cutoff vanished,
odd-sector vanishing would make the critical sum finite. This proves
the final assertion.
\end{proof}

The integral criterion~\eqref{eq:noiseintegral} is an exact positive
and negative statement for this specified correlated-clock coupling.
Theorem~\ref{thm:critical} is a negative statement for the original
clock-sector energies. We do not infer a positive subcritical
weighted-sector theorem, an equivalence with a completed fractional
clock-operator domain, or a physical correction topology. The parameter
$h$ changes Gaussian correlation, not physical observation time.
The Hermite correlation mechanism is classical Mehler theory;
Gaussian interpolation and fractional smoothness have extensive prior
theory, for example~\cite{geiss}. That literature is context rather
than a substitute for the particular source-valued coupling and domain
identifications proved here. A bounded preliminary search did not
settle priority for the exact reflected-clock energy application.

\section{Graph-dual derivatives and escape of the original field norm}
We give the closed-operator argument separately to
make its topology and representability exact.
Let $T:D(T)\subset E_0\to F_0$ be a closed densely
defined operator between real Hilbert spaces. Put
\[
 \mathcal G=D(T^*),\qquad
 \|v\|_{\mathcal G}^2=\|v\|_{F_0}^2+\|T^*v\|_{E_0}^2.
\]
The adjoint is closed, so this is a Hilbert space.
Define bounded maps
\begin{equation}\label{eq:dualdefinition}
 df(v)=\langle f,T^*v\rangle_{E_0},\qquad
 jV(v)=\langle V,v\rangle_{F_0}.
\end{equation}
Cauchy--Schwarz gives $\|df\|\leq\|f\|$ and
$\|jV\|\leq\|V\|$. For $u\in D(T)$,
\begin{equation}\label{eq:agreement}
 jTu=du.
\end{equation}

\begin{proposition}\label{prop:represent}
The map $j$ is injective, and $df=jV$ for some
$V\in F_0$ if and only if $f\in D(T)$, in which
case $V=Tf$. If $f\notin D(T)$, then $df\ne0$.
\end{proposition}
\begin{proof}
Write $\Gamma=\operatorname{graph}T$ in $E_0\oplus F_0$.
By the definition of the adjoint and density of $D(T)$,
\[
 \Gamma^\perp=\{(-T^*v,v):v\in D(T^*)\}.
\]
The identity $df=jV$ says exactly that $(f,V)$ is
orthogonal to this orthogonal complement. Since
$\Gamma$ is closed, $\Gamma^{\perp\perp}=\Gamma$;
thus $(f,V)$ belongs to the original graph. The
converse is \eqref{eq:agreement}.
If $jV=0$, the same orthogonality argument puts
$(0,V)$ in $\Gamma$, so $V=T0=0$. This also proves
that $D(T^*)$ is dense in $F_0$. Finally $df=0=j0$
would put $f$ in $D(T)$, proving the last assertion.
\end{proof}

For any $u\in D(T)$ and $f\in E_0$, the graph-dual
error bound is immediate:
\begin{equation}\label{eq:error}
 \|jTu-df\|_{\mathcal G^*}
 =\|d(u-f)\|_{\mathcal G^*}\leq\|u-f\|_{E_0}.
\end{equation}
Density of $D(T)$ supplies such approximating values
for every $f$. This convergence is in the operator norm
of the graph-test dual, not in the original field space.

\begin{proposition}\label{prop:escape}
If $u_n\in D(T)$, $u_n\to f$ in $E_0$, and
$f\notin D(T)$, then $\|Tu_n\|_{F_0}\to\infty$.
\end{proposition}
\begin{proof}
If the norms did not tend to infinity, some subsequence
would have $\|Tu_{n_j}\|\leq C$ for a finite $C$.
For $v\in D(T^*)$ take the limit in
$\langle u_{n_j},T^*v\rangle=\langle Tu_{n_j},v\rangle$:
\[
 |\langle f,T^*v\rangle|\leq C\|v\|_{F_0}.
\]
Since $D(T^*)$ is dense in $F_0$, this functional
extends continuously to all of $F_0$. Riesz
representation provides $V\in F_0$ with
$\langle f,T^*v\rangle=\langle V,v\rangle$ for
every adjoint-domain test. Proposition~\ref{prop:represent}
would put $f$ in $D(T)$, a contradiction.
\end{proof}

Apply these results to $T=D_W$, the closed densely
defined ambient clock lift of Section~\ref{sec:gaussian},
and $f=Z_T$ or $S_T$ (with the subscript $T>0$ now denoting time).
Theorems~\ref{thm:point} and~\ref{thm:bounded} give nonmembership, and
\eqref{eq:dualdefinition}, \eqref{eq:error} and
Proposition~\ref{prop:escape} prove
Theorem~\ref{thm:dualpoint}. The bound on the weak
derivative is explicitly
$\|d_WZ_T\|\leq(2T/\pi)^{1/4}$, while
$\|d_WS_T\|\leq1$.
Every original-domain value approximation has divergent
clock energy, although its embedded derivatives converge
in a weaker, precisely defined norm.

\section{Scope, antecedents and formal correspondence}
The Brownian-time model is classical~\cite{alloubazheng}.
Gaussian derivatives act on its independent sources;
$Z$ itself is not Gaussian. Standard derivative/adjoint
conventions appear in~\cite{sottinen}. The Brownian-time
change-of-measure work~\cite{change} addresses a different
question; none of the exclusions here rules out every
possible Brownian-time calculus. Novelty remains unresolved.

All formal correspondences use frozen verified commit
\begin{center}\small\ttfamily
bab688bb6140cd6316c80680095a914ea0d346db.
\end{center}
The accompanying theorem map records exact sources for
the folded-section metric, primitive and contraction
arguments, point and sine exclusions, clock-noise and
critical-sector energies, and generic graph-dual
representation/error/escape theorems. The ambient clock
application remains distinct from the separate native
adjoint-graph-dual time-average results.

The formal realization may contain unused independent
scenery coordinates, including a negative-time branch.
The half-line model retains the same process and partial
domains on its observables: graph pairs tensor with
constants in unused coordinates; conversely averaging
those coordinates contracts the graph test identities.
They do not provide a substitute clock gradient.

Value regularization alone cannot produce a convergent
original clock-gradient sequence at the point or sine;
closedness and the norm-escape theorem exclude it.
Scalar tests without one common finite-energy field do
not meet the recovered-field hypothesis. Neither a
source-variable divergence nor an adjoint-graph test
functional is identified as an observed-past physical-time
integrand. No physical-time It\^o/Clark--Ocone completion
or changing-past integrated point-regularizer limit is proved.

\begin{thebibliography}{9}
\bibitem{alloubazheng} H. Allouba and W. Zheng,
Brownian-time processes: the PDE connection and the
half-derivative generator, \emph{Ann. Probab.}
\textbf{29} (2001), 1780--1795.
\url{https://arxiv.org/abs/1005.3801}.
\bibitem{sottinen} T. Sottinen and L. Viitasaari,
Stochastic analysis of Gaussian processes via Fredholm
representation, \emph{Int. J. Stoch. Anal.}
\textbf{2016}, article 8694365.
\href{https://doi.org/10.1155/2016/8694365}
{doi:10.1155/2016/8694365}.
\bibitem{geiss} S. Geiss and A. Toivola,
On fractional smoothness and $L_p$-approximation on the Gaussian space,
\emph{Ann. Probab.} \textbf{43} (2015), 605--638.
\href{https://doi.org/10.1214/13-AOP884}{doi:10.1214/13-AOP884}.
\url{https://arxiv.org/abs/1206.5415}.
\bibitem{change} B. Abdurakhmanov and H. Allouba,
Brownian-time change of measure, preprint,
arXiv:2605.07044 (2026).
\url{https://arxiv.org/abs/2605.07044}.
\end{thebibliography}
\end{document}
